\part{Funzioni Operative}

\chapter{Signal K e Qualità dei Dati}
\label{ch:qualitadati}

In mare, un numero sullo schermo è utile solo se sai se è ancora fresco.
FairWindSK usa tre stati dei dati in modo uniforme in tutta l'interfaccia.

\section{Stati dei Singoli Valori}

\rowcolors{2}{LtGray}{white}
\begin{longtable}{L{2.4cm}L{5.2cm}L{6.1cm}}
\toprule
\rowcolor{Navy}
\color{white}\sffamily\bfseries Stato &
\color{white}\sffamily\bfseries Significato &
\color{white}\sffamily\bfseries Azione consigliata \\
\midrule
\textbf{Live}         &
  Valore ricevuto di recente e in aggiornamento normale. &
  Usa normalmente. Continua a guardare la Barra Superiore periodicamente. \\
\textbf{Non aggiornato} &
  Gli aggiornamenti si sono rallentati o fermati. Viene mostrato l'ultimo valore
  noto, attenuato. Un avviso, non la verità attuale. &
  Tratta con cautela. Controlla con strumenti indipendenti.
  Indaga se lo stato persiste. \\
\textbf{Mancante}     &
  Nessun valore utilizzabile. Viene mostrato un trattino o un segnaposto.
  \textbf{Mancante non significa zero.} Significa non disponibile. &
  Tratta come completamente sconosciuto.
  Non usare mai come valore di default sicuro.
  Verifica con un'altra fonte. \\
\bottomrule
\end{longtable}

\section{Stati di Salute della Connessione}

\rowcolors{2}{LtGray}{white}
\begin{longtable}{L{3.2cm}L{4.8cm}L{5.7cm}}
\toprule
\rowcolor{Navy}
\color{white}\sffamily\bfseries Stato &
\color{white}\sffamily\bfseries Significato &
\color{white}\sffamily\bfseries Risposta \\
\midrule
Live               & Completamente integra. &
  Continua normalmente. \\
Connessione        & Stabilendo la connessione iniziale. &
  Attendi 30--60 secondi. \\
Riconnessione      & Connessione persa; recupero automatico in corso. &
  Attendi brevemente. Di solito si risolve automaticamente. \\
Non aggiornato     & Connesso ma aggiornamenti rallentati o ritardati. &
  Controlla gli strumenti. Indaga se persiste. \\
REST degradato     & Canale richieste/risposte compromesso. &
  Le pagine Impostazioni potrebbero fallire; monitora. \\
Stream degradato   & Stream live compromesso. &
  Le letture potrebbero smettere di aggiornarsi; tratta come obsoleto. \\
Disconnesso        & Nessuna connessione. &
  Controlla il server Signal~K e il Wi-Fi. \\
App in caricamento & Catalogo app in recupero. &
  Attendi; le app appariranno quando pronto. \\
App non aggiornat. & Il catalogo app potrebbe non essere aggiornato. &
  Le app funzionano ancora; l'elenco potrebbe essere incompleto. \\
App primo piano deg.& L'applicazione corrente ha un problema. &
  Riprova una volta; se ancora in errore, torna al launcher. \\
\bottomrule
\end{longtable}

\section{Quando Agire Immediatamente}

Agisci senza indugio quando:

\begin{itemize}
  \item la profondità diventa obsoleta o mancante in acque basse o ristrette
  \item la posizione o la rotta scompaiono durante la navigazione in spazi ristretti
  \item si verifica un evento POB (Persona in Mare)
  \item un comando all'autopilota non corrisponde alla risposta reale dell'imbarcazione
\end{itemize}

Monitora brevemente e poi verifica quando:

\begin{itemize}
  \item la shell si sta riconnettendo dopo una breve interruzione
  \item il caricamento delle app richiede più tempo del solito
  \item i valori di guida rotta smettono di aggiornarsi
  \item solo una app sembra degradata mentre il resto della shell è integro
\end{itemize}

\section{Ciclo di Recupero Automatico}

Quando viene persa una connessione, FairWindSK automaticamente:

\begin{enumerate}
  \item Rileva la disconnessione o il degrado.
  \item Esegue un nuovo rilevamento del server Signal~K.
  \item Ristabilisce lo stream WebSocket.
  \item Reinvia tutte le sottoscrizioni attive.
  \item Esegue uno snapshot REST immediato affinché i widget mostrino valori aggiornati.
  \item Aggiorna il catalogo app e tutti i modelli di risorse MyData.
  \item Aggiorna le preferenze di unità dal server.
\end{enumerate}

\begin{warningbox}
Dopo un'interruzione di corrente, i ricevitori GPS e gli strumenti impiegano tempo
per riacquisire e stabilizzarsi.
Anche dopo che FairWindSK mostra \textit{Live}, verifica ogni valore critico per
la sicurezza prima di affidarti ad esso.
\end{warningbox}

\chapter{Persona in Mare (POB)}
\label{ch:pob}

\begin{warningbox}
Segnala immediatamente un evento di persona in mare.
Non ritardare perché pensi di ricordare la posizione dopo.
Segnala prima, poi continua le azioni di recupero.
La funzione POB assiste nel riferimento per il recupero.
Non sostituisce la risposta immediata dell'equipaggio, la guardia visiva,
il timone o le procedure di emergenza consolidate.
\end{warningbox}

\section{Attivazione POB}

Tocca \iconlabel{alarm-pob}{POB} nella Barra Inferiore.
È sempre visibile e mai disabilitato.
Esercitati a trovarlo al buio e in condizioni difficili.

Nell'istante in cui tocchi POB, FairWindSK simultaneamente:

\begin{itemize}
  \item Acquisisce la posizione GPS corrente come posizione dell'incidente.
  \item Crea una risorsa waypoint su Signal~K con timestamp, denominata
        \code{POB HH:MM:SS} (UTC).
  \item Imposta quel waypoint come destinazione di navigazione attiva ---
        tutte le app di navigazione mostrano immediatamente rilevamento e distanza dal POB.
  \item Attiva l'allarme standard Signal~K \skpath{notifications.mob} a livello di emergenza,
        avvisando tutti i dispositivi connessi.
  \item Avvia il cronometro del tempo trascorso, che conta ogni secondo.
  \item Si abbona a rilevamento e distanza in tempo reale dall'imbarcazione al waypoint POB.
  \item Sostituisce la Barra Inferiore normale con il dialogo POB.
\end{itemize}

\screenshot{pob-bar-active}{Barra POB: rilevamento, distanza, tempo trascorso, pulsanti ricerca SAR}{}

\section{Controlli e Letture della Barra POB}

\rowcolors{2}{LtGray}{white}
\begin{longtable}{C{1.4cm}L{3.0cm}L{9.3cm}}
\toprule
\rowcolor{Navy}
\color{white}\sffamily\bfseries Icona &
\color{white}\sffamily\bfseries Controllo &
\color{white}\sffamily\bfseries Descrizione \\
\midrule
---                      & Selettore POB  &
  Elenco a discesa di tutti gli eventi POB attivi.
  Seleziona un evento per visualizzarne i dati. \\
---                      & Ultima pos. nota &
  Lat e lon al momento della pressione di POB, nel formato coordinate configurato. \\
---                      & Rilevamento     &
  Rilevamento bussola in tempo reale dall'imbarcazione al waypoint POB.
  Si aggiorna continuamente. \\
---                      & Distanza        &
  Distanza in tempo reale dall'imbarcazione al waypoint POB. Si aggiorna continuamente. \\
---                      & Tempo trascorso &
  HH:MM:SS che conta dal momento della pressione di POB. \\
---                      & Annulla         &
  Annulla l'incidente selezionato.
  Rimuove l'allarme Signal~K e la destinazione di navigazione per quell'evento. \\
\iconlg{search-spiral}   & Ricerca a spirale &
  Genera un percorso SAR a spirale centrato sulla posizione POB. \\
\iconlg{search-parallel} & Ricerca parallela &
  Genera un percorso SAR a linee parallele centrato sulla posizione POB. \\
\iconlg{close-google}    & Chiudi dialogo  &
  Torna alla Barra Inferiore normale.
  \textbf{L'allarme e il waypoint rimangono attivi su Signal~K.} \\
\bottomrule
\end{longtable}

\begin{warningbox}
Toccare \iconlg{close-google}~Chiudi \textbf{non} annulla il POB.
L'allarme Signal~K e il waypoint rimangono attivi finché non vengono esplicitamente annullati.
Solo il pulsante \key{Annulla} rimuove un incidente specifico.
\end{warningbox}

\section{Schemi di Ricerca SAR}

\subsection{Ricerca a Spirale}

Tocca \iconlabel{search-spiral}{Ricerca a spirale}.
FairWindSK crea una regione SAR di $500 \times 500$\,m centrata sulla posizione POB
e un percorso a spirale che parte dal POB e si irradia verso l'esterno con spaziatura di 50\,m.
Sia la regione che il percorso vengono salvati come risorse Signal~K e appaiono immediatamente
in \ui{MyData~> Regioni} e \ui{MyData~> Rotte}.

\subsection{Ricerca a Linee Parallele}

Tocca \iconlabel{search-parallel}{Ricerca parallela}.
FairWindSK crea la stessa regione SAR di $500 \times 500$\,m e un percorso a linee
parallele che la copre con spaziatura di 60\,m.

\section{Annullare un POB}

\begin{enumerate}
  \item Seleziona l'incidente nel selettore POB (se ci sono più incidenti attivi).
  \item Tocca \key{Annulla}.
  \item FairWindSK rimuove l'allarme Signal~K e cancella la destinazione di navigazione.
  \item Se non ci sono più incidenti attivi, il dialogo POB si chiude automaticamente.
\end{enumerate}

\begin{notebox}
Il waypoint POB \textbf{non} viene eliminato all'annullamento.
Viene conservato in \ui{MyData~> Waypoint} con il suo timestamp e la sua posizione
per la documentazione di eventuali incidenti e il debriefing post-incidente.
\end{notebox}

\chapter{Controlli Autopilota}
\label{ch:autopilota}

Tocca \iconlabel{command-autopilot}{Autopilota} nella Barra Inferiore per aprire la barra
Autopilota, che sostituisce la Barra Inferiore normale.
L'icona Autopilota è attiva solo quando il plugin richiesto è installato e FairWindSK
può interagire con esso.

\begin{warningbox}
Verifica sempre l'effetto reale di un comando autopilota sull'imbarcazione.
Non dare mai per scontato che un cambio di modalità sia avvenuto solo perché hai toccato
un controllo.
Sii pronto a disinnestare o sovrascrivere immediatamente se la risposta è pericolosa
o inattesa.
\end{warningbox}

\screenshot{autopilot-bar}{Barra Autopilota: pulsanti modalità, controlli di assetto, indicatore timone}{}

\section{Controlli della Barra Autopilota}

Le icone sono gli asset esatti definiti in \code{AutopilotBar.cpp}.

\rowcolors{2}{LtGray}{white}
\begin{longtable}{C{1.4cm}L{2.4cm}L{4.6cm}L{5.3cm}}
\toprule
\rowcolor{Navy}
\color{white}\sffamily\bfseries Icona &
\color{white}\sffamily\bfseries Controllo &
\color{white}\sffamily\bfseries Azione Signal~K &
\color{white}\sffamily\bfseries Quando usarlo \\
\midrule
\iconlg{close-google}                & Standby      & Disinnesta           & Qualsiasi volta si voglia il timone manuale. \\
\iconlg{command-autopilot}           & Auto         & Modalità bussola     & Traversate in mare aperto; mantiene la prua. \\
\iconlg{up-iec}                      & Vento        & Modalità vento       & Bolina/poppa; mantiene l'angolo di vento. \\
\iconlg{navigation-route}            & Rotta        & Modalità GPS         & Navigazione su percorso waypoint. \\
\iconlg{spinbox-minus}               & $-1\textdegree$ sx & Assetto $-1\textdegree$ & Regolazione fine a sinistra. \\
\iconlg{spinbox-plus}                & $+1\textdegree$ dr & Assetto $+1\textdegree$ & Regolazione fine a dritta. \\
\iconlg{chevron-double-left-google}  & $-10\textdegree$ sx & Correggi $-10\textdegree$ & Correzione di rotta maggiore a sinistra. \\
\iconlg{chevron-double-right-google} & $+10\textdegree$ dr & Correggi $+10\textdegree$ & Correzione di rotta maggiore a dritta. \\
\iconlg{arrow-left-google}           & Virata sin.  & Esegui virata sx     & Navigazione a vela in bolina. \\
\iconlg{arrow-right-google}          & Virata dr.   & Esegui virata dr     & Navigazione a vela in bolina. \\
\iconlg{chevron-double-left-google}  & Strambata sx & Esegui strambata sx  & Condizioni da lievi a moderate; supervisionare. \\
\iconlg{chevron-double-right-google} & Strambata dr & Esegui strambata dr  & Supervisionare al timone in vento forte. \\
\iconlg{waypoint-optional-iec}       & WPT succ.    & Avanza leg rotta     & Navigazione su rotta. \\
\iconlg{alerts}                      & Scarto       & Manovra evitamento   & Nassa, imbarcazione, ostacolo. \\
\bottomrule
\end{longtable}

\section{Indicatore Timone}

La barra Autopilota mostra uno slider orizzontale \textbf{di sola lettura} che visualizza
la deflessione corrente del timone da $-30\textdegree$ (tutto a sinistra) a
$+30\textdegree$ (tutto a dritta).
Riflette la posizione effettiva del timone riportata dal pilota --- non una prua comandata.
Usalo per confermare che l'autopilota risponde ai comandi come atteso.

\chapter{Guardia Ancora}
\label{ch:ancora}

Tocca \iconlabel{anchor-iec}{Ancora} per aprire la barra Guardia Ancora.

\begin{warningbox}
La guardia ancora elettronica è solo un ausilio.
In un ancoraggio esposto, con tempo in peggioramento o con fondale scarso,
mantieni una guardia umana indipendentemente da ciò che mostra qualsiasi
dispositivo elettronico.
\end{warningbox}

\screenshot{anchor-bar}{Barra Guardia Ancora: posizione, profondità, calumo, raggio, controlli argano}{}

\section{Controlli e Letture della Barra Ancora}

Le icone sono gli asset esatti definiti in \code{AnchorBar.ui}.

\rowcolors{2}{LtGray}{white}
\begin{longtable}{C{1.4cm}L{3.0cm}L{9.3cm}}
\toprule
\rowcolor{Navy}
\color{white}\sffamily\bfseries Icona &
\color{white}\sffamily\bfseries Controllo &
\color{white}\sffamily\bfseries Descrizione \\
\midrule
\iconlg{anchor-drop}    & Cala          & Registra la posizione GPS come punto di ancoraggio e avvia il monitoraggio. \\
\iconlg{anchor-raise}   & Salpa         & Ferma il monitoraggio e cancella la posizione dell'ancora. \\
\iconlg{anchor-release} & Fila          & Invia comando di fila catena all'argano elettrico (se supportato). \\
\iconlg{anchor-up}      & Su (tieni)    & Comando continuo di alaggio all'argano elettrico. Rilascia per fermare. \\
\iconlg{anchor-down}    & Giù (tieni)   & Comando continuo di filatura all'argano elettrico. Rilascia per fermare. \\
\iconlg{anchor-reset}   & Reimposta     & Reimposta la posizione dell'ancora alla posizione GPS corrente. \\
---                      & Slider raggio & Distanza massima consentita prima dell'allarme. \\
---                      & Imposta Raggio & Conferma il raggio e lo invia al plugin ancora. \\
---                      & LAT / LON     & Coordinate GPS del punto di ancoraggio registrato. \\
---                      & Profondità    & Profondità attuale dell'acqua dall'ecoscandaglio. \\
---                      & Calumo        & Lunghezza di calumo filato. \\
---                      & Raggio att./Rilev. & Distanza e rilevamento in tempo reale dall'imbarcazione all'ancora. \\
---                      & Raggio max    & Distanza massima raggiunta dall'ancora dall'inizio della guardia. \\
\bottomrule
\end{longtable}

\section{Impostazione della Guardia Ancora --- Passo per Passo}

\begin{enumerate}
  \item Ancora normalmente e fila il calumo previsto.
  \item Dopo che l'imbarcazione si è assestata, tocca \iconlg{anchor-iec} nella Barra Inferiore.
  \item Conferma che la posizione GPS visualizzata corrisponda a dove giace l'ancora.
  \item Tocca \iconlg{anchor-drop}~\key{Cala}. FairWindSK registra la posizione e avvia il monitoraggio.
  \item Imposta il raggio, tenendo conto di: calumo, cerchio di bordata, precisione GPS
        ($\pm3$--$10$\,m) e margine di sicurezza.
  \item Tocca \key{Imposta Raggio} per confermare.
  \item Chiudi la barra. Il monitoraggio continua in background.
\end{enumerate}

\chapter{Allarmi}
\label{ch:allarmi}

Tocca \iconlabel{alarm-general}{Allarmi} per aprire la barra Allarmi.
Quando un allarme è attivo, l'icona viene evidenziata con il colore accent su tutti i preset
comfort --- progettato per attirare l'attenzione anche durante la guardia notturna.

\screenshot{alarms-bar}{Barra Allarmi: pulsanti di emergenza SOLAS ed elenco allarmi attivi}{}

\section{Pulsanti di Allarme di Emergenza}

Le icone sono gli asset esatti definiti in \code{AlarmsBar.ui}.

\rowcolors{2}{LtGray}{white}
\begin{longtable}{C{1.4cm}L{2.2cm}L{4.0cm}L{6.5cm}}
\toprule
\rowcolor{Navy}
\color{white}\sffamily\bfseries Icona &
\color{white}\sffamily\bfseries Pulsante &
\color{white}\sffamily\bfseries Percorso Signal~K &
\color{white}\sffamily\bfseries Tipo di emergenza \\
\midrule
\iconlg{alarm-pob}     & POB          & \skpath{notifications.mob}     & Persona / uomo in mare \\
\iconlg{alarm-sinking} & Affondamento & \skpath{notifications.sinking} & Imbarcazione che imbarca acqua o affonda \\
\iconlg{alarm-fire}    & Incendio     & \skpath{notifications.fire}    & Incendio a bordo \\
\iconlg{alarm-piracy}  & Pirateria    & \skpath{notifications.piracy}  & Pirateria o attacco armato \\
\iconlg{alarm-abandon} & Abbandono    & \skpath{notifications.abandon} & Abbandono della nave \\
\iconlg{alarm-adrift}  & Alla deriva  & \skpath{notifications.adrift}  & Imbarcazione alla deriva e incontrollata \\
\bottomrule
\end{longtable}

\begin{warningbox}
L'attivazione di un allarme in FairWindSK genera una notifica Signal~K solo sulla
rete di bordo.
\textbf{Non} contatta la guardia costiera, l'MRCC o alcuna autorità esterna.
Effettua una chiamata di soccorso DSC e/o vocale sul Canale VHF~16 in qualsiasi
emergenza reale.
\end{warningbox}

\section{Priorità degli Allarmi e Risposta}

Gestisci gli allarmi in questo ordine:

\begin{enumerate}
  \item Prima i valori critici per la sicurezza --- profondità, posizione, rotta.
  \item Poi la salute della connessione e della shell.
  \item Infine i problemi specifici dell'applicazione.
\end{enumerate}

\section{Elenco Allarmi Attivi}

Sotto i pulsanti, la barra Allarmi elenca tutte le notifiche Signal~K attualmente attive,
mostrando percorso, stato (emergenza/allarme/avviso), testo del messaggio e timestamp.
Mostra anche lo stato corrente della connessione Signal~K e il timestamp dell'ultimo
aggiornamento dello stream --- un sommario compatto della salute senza dover uscire dalla barra.
